В рамках выполнения выпускной квалификационной работы была разработана библиотека криптографических примитивов Curve25519 и Ed25519 для платформы .NET.

Полный исходный код библиотеки, а также инструкции по сборке и примеры использования размещены в открытом репозитории на платформе GitHub.


\textbf{Ссылка на репозиторий:} \url{https://github.com/yashelter/dotnet-curve25519-ed25519}


\begin{figure}[H]
    \centering
    \qrcode[height=5cm]{https://github.com/yashelter/dotnet-curve25519-ed25519}
    
    \vspace{1ex}
    \caption{QR-код со ссылкой на репозиторий с исходным кодом реализации}
    \label{fig:qr_code_github}
\end{figure}

В состав репозитория входят:


\begin{enumerate}
    \item Исходный код вычислительного ядра (Radix-51, AVX2 Batching).
    \item Реализация алгоритмов X25519 и Ed25519.
    \item Проекты для проведения нагрузочного тестирования (BenchmarkDotNet).
    \item Набор юнит-тестов, проверяющих соответствие эталонным векторам RFC.
\end{enumerate}